(1) 入射总能量为 \(E=(\gamma m_1+m_2)c^2\)，总动量 \(p=\gamma m_1v\)。合体能量与动量分别为 \(\Gamma Mc^2,\Gamma MV\)，两者之比给 \(V=pc^2/E=\gamma m_1v/(\gamma m_1+m_2)\)，满足 \(0<V<v<c\)。
(2) 四动量平方守恒 \(M^2c^4=E^2-p^2c^2\)；将 \(\gamma^2(1-v^2/c^2)=1\) 化简得 \(M=\sqrt{m_1^2+m_2^2+2\gamma m_1m_2}\)。当 \(v>0\)，\(\gamma>1\)，于是 \(M^2>(m_1+m_2)^2\)，故 \(M>m_1+m_2\)，增加的是合体内部能的质量当量。
(3) 静质量增量对应内部能 \(\Delta U_{\rm int}=[M-(m_1+m_2)]c^2>0\)。实验室系能量仍严格守恒：合体总能 \(\Gamma Mc^2=(\gamma m_1+m_2)c^2\)，动量 \(\Gamma MV=\gamma m_1v\)；多出的静能来自被损失的相对运动动能而非外界额外输入。
(4) 对 \(v/c\ll1\)，\(\gamma=1+v^2/(2c^2)+O(v^4/c^4)\)，于是 \(V=\frac{m_1}{m_1+m_2}v+O(v^3/c^2)\)。再展开 \(M=(m_1+m_2)+[m_1m_2/(2(m_1+m_2))](v^2/c^2)+O(v^4/c^4)\)，故 \(\Delta U_{\rm int}=\frac12\mu v^2+O(v^4/c^2)\)，\(\mu=m_1m_2/(m_1+m_2)\)，正是经典完全非弹性碰撞的相对动能损失。
